% A tabela \ref{tab:limites_juntasPosVelocidade_rad} a seguir complementa a tabela \eqref{tab:limites_juntasPosVelocidade} com os dados convertidos, no caso das juntas revolutas para radianos (rad) e radianos por segundo (rad/s), e para a junta prismática permanecendo em metro (m) e metro por segundo (m/s).

% \begin{table}[ht]
% \centering
% {
% \setlength{\tabcolsep}{12pt}
% \renewcommand{\arraystretch}{1}
% \caption{Limites físicos e operacionais das juntas do manipulador robótico}
% \label{tab:limites_juntasPosVelocidade_rad}
% \begin{tabular}{clrrrr}
% \toprule
% \textbf{Junta} & \textbf{Tipo} & $q_{j,\min}$ & $q_{j,\max}$ & $\dot{q}_{j,\min}$ & $\dot{q}_{j,\max}$ \\
% \midrule
% 1 & Revoluta   & $-6{,}2832$ & $6{,}2832$ &
% $-0{,}5\pi$ & $0{,}5\pi$ \\[4pt]
% 2 & Revoluta   & $-1{,}5708$ & $1{,}5708$ &
% $-\pi$ & $\pi$ \\[4pt]
% 3 & Revoluta   & $-2{,}2689$ & $2{,}2689$
% & $-\pi$ & $\pi$ \\[4pt]
% 4 & Revoluta   & $-2{,}6180$ & $2{,}6180$ &
% $-2\pi$ & $2\pi$ \\[4pt]
% 5 & Prismática & $-0{,}15$ m & $0{,}00$ m &
% $-0{,}314$ m/s & $0{,}314$ m/s \\[4pt]
% \bottomrule
% \end{tabular}
% }
% \end{table}

\chapter{Parâmetros Inerciais e Agregação Dinâmic}
\label{ApendiceC}
Após a modelagem geométrica do manipulador robótico em 3D utilizando a ferramenta CAD (\emph{Computer-Aided Design}) Autodesk Fusion 360, obtiveram-se os parâmetros inerciais individuais de cada elo, junta e base, bem como os parâmetros do manipulador completo.
Esses dados compreendem as massas, a posição do centro de massa local e o tensor de inércia em relação ao respectivo centro de massa e \emph{frame} local. Todos os valores detalhados encontram-se no Apêndice \eqref{ap:ApendiceA}.

\subsection{Parâmetros inerciais por elo}
Como o centro de massa é fornecido no sistema de referência local do elo, é necessário convertê-lo para o sistema de referência da base $\{0\}$.
Sejam ${}^{0}\!\mathbf R_i(q)\in\mathbb{R}^{3\times 3}$ a rotação do frame do corpo $i$ em relação à base e ${}^{0}\!\vec p_i(q)$ a posição de sua origem. O centro de massa do corpo $i$ expresso na base é
\begin{equation}
{}^{0}\!\vec r_{C_i}(q)= {}^{0}\!\vec p_i(q)+{}^{0}\!\mathbf R_i(q)\,{}^{i}\!\vec r_{C_i}.
\label{eq:rci_base}
\end{equation}

\subsection{Transformação para o frame da base}
De forma análoga, o tensor de inércia ${}^{(elo)}\mathbf{I}_{i}$, o tensor de inércia do corpo ${}^{(elo)}\mathbf{I}_{i}$, quando expresso no seu centro de massa, e alinhado aos eixos principais locais, é rotacionado para os eixos da base por:
\begin{equation}
{}^{0}\!\mathbf I_{C_i}(q)= {}^{0}\!\mathbf R_i(q)\,{}^{(i)}\!\mathbf I_{C_i}\,{}^{0}\!\mathbf R_i^{\top}(q).
\label{eq:I_rot_base}
\end{equation}

\subsection{Aplicação do Teorema de Steiner}
Mesmo após a rotação, o tensor de inércia ainda está referido ao centro de massa do elo, para que esteja em referência a inércia do corpo $i$ a um ponto $p$ expresso na base, utiliza-se o Teorema de Steiner \eqref{eq:steiner_matriz} para referi-lo à origem da base $O$:
\begin{equation}
{}^{0}\!\mathbf I_{P_i}(q)
= {}^{0}\!\mathbf I_{C_i}(q)+ m_i\big(\,\|{}^{0}\!\vec d_{Pi}(q)\|^2\mathbf I_3-{}^{0}\!\vec d_{Pi}(q)\,{}^{0}\!\vec d_{Pi}^{\,T}(q)\big),
\label{eq:steiner_matriz}
\end{equation}

Onde ${}^{0}\mathbf{I}_{O_i}$ é o tensor de inércia do elo $i$ em relação à origem da base;
${}^{0}\mathbf{I}_{C_i}$ é o tensor no centro de massa já rotacionado para a base;
$m_i$ é a massa do elo $i$;
$\vec r_i$ é a posição do seu centro de massa em relação à base;
$\mathbf{I}_3$ é a matriz identidade $3\times 3$ e por fim,
${}^{0}\!\vec d_{Pi}(q)= {}^{0}\!\vec r_{C_i}(q)-{}^{0}\!\vec r_{P}$.
%% ===== EXPLICAÇÃO =====
% A matriz identidade serve para construir um tensor de inercia equivalente que leve em conta a mudança de referencia (do CoM para outro ponto)
% I_3 aparece pq estamos tratando grandezas vetoriais 3D
% é a forma matricial de generalizar a formula escalar I = sum mr^2 para o caso tensorial
% sem ela, não teriamos a componente isotropica


\subsection{Agregação dos parâmetros}
A seguir todos os parâmetros já calculados serão agregados, iniciando pela massa.
A massa total do conjunto (base + manipulador robótico) é dado por:
\begin{equation}
  M_{\text{total}} = \sum_{j} m_j
\end{equation}

O centro de massa global, expresso na base, é:
\begin{equation}
  \vec{r}_{\text{Global}} = \frac{1}{M_{\text{total}}} \sum_{j} m_j \,\vec{r}_j .
\end{equation}

Para encontrar a inércia total do sistema referida à origem da base:
\begin{equation}
{}^{0}\!\mathbf I_{O}(q)=\sum_{i} {}^{0}\!\mathbf I_{O_i}(q)
=\sum_{i}\Big[\,{}^{0}\!\mathbf I_{C_i}(q)+ m_i\big(\|{}^{0}\!\vec r_{C_i}(q)\|^2\mathbf I_3-{}^{0}\!\vec r_{C_i}(q)\,{}^{0}\!\vec r_{C_i}^{\,T}(q)\big)\Big].
\label{eq:I_total_O}
\end{equation}

A inércia equivalente no centro de massa global é:
\begin{equation}
{}^{0}\!\mathbf I_{G}(q)
= {}^{0}\!\mathbf I_{O}(q)
- M_{\mathrm{tot}}\big(\|{}^{0}\!\vec r_{G}(q)\|^2\mathbf I_3-{}^{0}\!\vec r_{G}(q)\,{}^{0}\!\vec r_{G}^{\,T}(q)\big).
\label{eq:I_total_G}
\end{equation}

\subsection{Observações e verificação numérica}
Como ${}^{0}\!\mathbf R_i(q)$ e ${}^{0}\!\vec p_i(q)$ dependem de $q(t)$, as matrizes \({}^{0}\!\mathbf I_{O}(q)\) e \({}^{0}\!\mathbf I_{G}(q)\) variam ao passar do tempo.
Pode-se utilizar a propriedade a seguir para evitar assimetrias numéricas por arredondamento:
\begin{equation}
\mathbf I \;\leftarrow\; \tfrac{1}{2}\big(\mathbf I+\mathbf I^{T}\big).
\label{eq:symmetrize}
\end{equation}

As grandezas agregadas
\({}^{0}\!\mathbf I_{O}(q)\),
\({}^{0}\!\mathbf I_{G}(q)\) e
\({}^{0}\!\vec r_{G}(q)\)
alimentam diretamente a expressão de \(T(q,\dot q)\)
e a construção de \(\mathbf M(q)\) na seção \eqref{sec:energia_cinetica};
a formulação Lagrangiana e a forma \(\mathbf M(q)\ddot q+\mathbf C(q,\dot q)\dot q+\mathbf g(q)=\boldsymbol\tau\) estão na seção \eqref{sec:equacao_euler_lagrange}. 

\subsection{Parâmetros físicos do manipulador robótico}
A seguir são apresentados os parâmetros físicos do manipulador robótico, incluindo massas, dimensões e matrizes de inércia. 
Os valores encontram-se no Sistema Internacional (SI), com intuito de manter a consistência com a modelagem matemática desenvolvida. 
Para maior detalhamento, as mesmas tabelas também estão disponíveis em outras unidades (g e cm) no Apêndice
\eqref{ap:ApendiceA},
onde é possível verificar todas as casas decimais e observar que os produtos de inércia, embora muito pequenos, não são nulos.

\subsubsection{Massas e dimensões}
As massas e dimensões físicas de cada componente do manipulador robótico foram medidas a partir do modelo tridimensional no \emph{Autodesk Fusion 360} e convertidas para o Sistema Internacional. 
Esses parâmetros foram utilizados para a formulação da modelagem geométrica e para a determinação das matrizes de inércia, sendo posteriormente utilizados nas simulações computacionais. 
A Tabela \eqref{tab:dimensoes_massas_m_kg} apresenta o comprimento, o raio, o diâmetro e a massa de cada componente, bem como o valor total da massa do manipulador robótico.

\begin{table}[H]
\centering
\caption{Dimensões e massas dos componentes do manipulador (Sistema Internacional).}
\label{tab:dimensoes_massas_m_kg}
\begin{tabular}{l rrrr}
\hline
\textbf{Componente} & \textbf{Comprimento [m]} & \textbf{Raio [m]} & \textbf{Diâmetro [m]} & \textbf{Massa [kg]} \\
\hline
Base Geral   & 0{,}100 & 0{,}250  & 0{,}500 & 3{,}6530  \\
Junta 1      & 0{,}005 & 0{,}0250 & 0{,}050 & 0{,}0265 \\
Junta 2      & 0{,}005 & 0{,}0250 & 0{,}050 & 0{,}0265 \\
Junta 3      & 0{,}005 & 0{,}0250 & 0{,}050 & 0{,}0265 \\
Junta 4      & 0{,}005 & 0{,}0250 & 0{,}050 & 0{,}0265 \\
Junta 5      & 0{,}005 & 0{,}0250 & 0{,}050 & 0{,}0265 \\
Elo 1        & 0{,}100 & 0{,}0250 & 0{,}050 & 0{,}1115 \\
Elo 2        & 0{,}300 & 0{,}0250 & 0{,}050 & 0{,}3129 \\
Elo 3        & 0{,}300 & 0{,}0250 & 0{,}050 & 0{,}3129 \\
Elo 4        & 0{,}300 & 0{,}0250 & 0{,}050 & 0{,}3129 \\
Ferramenta   & 0{,}150 & 0{,}0075 & 0{,}015 & 0{,}0403 \\
\hline
\textbf{Total} & -- & -- & -- & \textbf{4{,}876} \\
\hline
\end{tabular}
\end{table}


\subsubsection{Momentos de inércia}
Os momentos principais de inércia de cada componente do manipulador foram obtidos através do software \emph{Fusion 360} e estão apresentados na Tabela \eqref{tab:inercias_diag_kgm2}. 
Os valores correspondem às componentes diagonais das matrizes de inércia ($I_{xx}$, $I_{yy}$ e $I_{zz}$), expressas em kg$\cdot$m$^{2}$. 
Apesar das magnitudes reduzidas, esses parâmetros são relevantes para a análise dinâmica, principalmente em sistemas espaciais onde pequenas variações podem impactar o desempenho do controle.

\begin{table}[H]
\centering
\caption{Momentos de inércia principais (valores diagonais) em kg$\cdot$m$^{2}$.}
\label{tab:inercias_diag_kgm2}
\begin{tabular}{l rrr}
\hline
\textbf{Componente} & $I_{xx}$ & $I_{yy}$ & $I_{zz}$ \\
\hline
Base Geral   & 0,07950   & 0,0795   & 0,144900 \\
Junta 1      & 4,20$\times$10$^{-6}$ & 4,20$\times$10$^{-6}$ & 8,28$\times$10$^{-6}$ \\
Junta 2      & 8,28$\times$10$^{-6}$ & 4,20$\times$10$^{-6}$ & 4,20$\times$10$^{-6}$ \\
Junta 3      & 8,28$\times$10$^{-6}$ & 4,20$\times$10$^{-6}$ & 4,20$\times$10$^{-6}$ \\
Junta 4      & 8,28$\times$10$^{-6}$ & 4,20$\times$10$^{-6}$ & 4,20$\times$10$^{-6}$ \\
Junta 5      & 4,20$\times$10$^{-6}$ & 8,28$\times$10$^{-6}$ & 4,20$\times$10$^{-6}$ \\
Elo 1        & 0,00137  & 0,00137  & 0,000597 \\
Elo 2        & 0,02580   & 0,00174  & 0,025800 \\
Elo 3        & 0,02580   & 0,00174  & 0,025800 \\
Elo 4        & 0,02580   & 0,00174  & 0,025800 \\
Ferramenta   & 7,82$\times$10$^{-5}$ & 1,62$\times$10$^{-6}$ & 7,82$\times$10$^{-5}$ \\
\hline
\end{tabular}
\end{table}



\subsubsection{Produtos de inércia}
Além dos momentos principais, os produtos de inércia ($I_{xy}$, $I_{xz}$ e $I_{yz}$) também foram mantidos e encontram-se listados na Tabela
\eqref{tab:inercias_prod_kgm2}. 
Embora apresentem valores de ordem muito pequena, ao considerar mais casas decimais permite manter a consistência matemática dos modelos e evita simplificações que poderiam comprometer análises de precisão. 
Essa escolha reforça a adequação do modelo a cenários de microgravidade, nos quais a dinâmica é sensível a perturbações de baixa intensidade.



\begin{table}[H]
\centering
\caption{Produtos de inércia (valores fora da diagonal) em kg$\cdot$m$^{2}$.}
\label{tab:inercias_prod_kgm2}
\begin{tabular}{l rrr}
\hline
\textbf{Componente} & $I_{xy}$ & $I_{xz}$ & $I_{yz}$ \\
\hline
Base Geral   & 6,60$\times$10$^{-16}$ & 1,64$\times$10$^{-12}$ & -1,30$\times$10$^{-16}$ \\
Junta 1      & -2,25$\times$10$^{-16}$ & 0 & 0 \\
Junta 2      & -1,13$\times$10$^{-17}$ & 0 & 0 \\
Junta 3      & -1,42$\times$10$^{-16}$ & 1,83$\times$10$^{-16}$ & -1,25$\times$10$^{-15}$ \\
Junta 4      & 3,08$\times$10$^{-16}$ & -1,32$\times$10$^{-17}$ & -4,35$\times$10$^{-16}$ \\
Junta 5      & -3,95$\times$10$^{-16}$ & 5,37$\times$10$^{-16}$ & -1,63$\times$10$^{-17}$ \\
Elo 1        & 0 & -1,17$\times$10$^{-14}$ & 2,35$\times$10$^{-17}$ \\
Elo 2        & 2,14$\times$10$^{-15}$ & -1,08$\times$10$^{-15}$ & 2,39$\times$10$^{-14}$ \\
Elo 3        & 9,95$\times$10$^{-17}$ & -6,59$\times$10$^{-15}$ & 4,41$\times$10$^{-14}$ \\
Elo 4        & 6,33$\times$10$^{-14}$ & 4,29$\times$10$^{-15}$ & -1,15$\times$10$^{-15}$ \\
Ferramenta   & 5,43$\times$10$^{-15}$ & 4,39$\times$10$^{-16}$ & -5,66$\times$10$^{-16}$ \\
\hline
\end{tabular}
\end{table}





